\section{Diskussion und Reflexion}
\label{sec:diskussion}

Die Diskussion der Ergebnisse zeigt, dass die Prognose des Aluminiumpreises eine herausfordernde Aufgabe ist, die durch hohe Volatilität und die Vielzahl von Einflussfaktoren geprägt wird.
Die entwickelten Systeme sind auf Basis der gewählten Daten und Modelle nicht in der Lage, eine Preisentwicklung mit einer signifikant höheren Wahrscheinlichkeit als $\approx 50\,\%$ vorherzusagen.

\subsection{Markteffizienz und Vorhersagbarkeit}
Das geringe Bestimmtheitsmaß ($R^2 \approx 0{,}0043$) und eine Richtungsgenauigkeit von $51{,}1\,\%$ stehen im Einklang mit der \emph{Efficient Market Hypothesis} (EMH).
Nach dieser Theorie spiegeln Marktpreise die öffentlich verfügbaren Informationen weitgehend wider, sodass Preisänderungen vor allem auf neue, nicht antizipierbare Ereignisse reagieren und damit nur begrenzt aus historischen Daten ableitbar sind.
Die Ergebnisse legen nahe, dass historische Preis- und Volumendaten allein nicht ausreichen, um belastbare Vorhersagemuster für den Aluminiummarkt zu extrahieren.
Dies dürfte auf die hohe Komplexität des Marktes zurückzuführen sein, in dem zahlreiche Faktoren wie geopolitische Ereignisse, Energiepreise oder technologische Entwicklungen eine Rolle spielen, die in den verwendeten Zeitreihen nur indirekt abgebildet werden.

\subsection{Bedeutung der probabilistischen Prognose und Unsicherheitskalibrierung}
Durch den Wechsel auf probabilistische Prognosen mittels Monte-Carlo-Dropout (MC-Dropout) kann das System seine Unsicherheit in Form von Konfidenzintervallen ausdrücken. Ein Konfidenzintervall ist ein Bereich, der mit einer bestimmten Wahrscheinlichkeit den tatsächlichen Preis enthalten soll. Die Ergebnisse in Tabelle~\ref{tab:results_lstm} zeigen dabei sowohl Stärken als auch Grenzen des Ansatzes. Bei $h=1$ liegt die Coverage bei $90{,}8\,\%$, was auf eine noch brauchbare Kalibrierung für sehr kurzfristige Vorhersagen hindeutet. Bei längeren Horizonten sinkt die Coverage jedoch deutlich ($73{,}2\,\%$ bei $h=5$ und $56{,}0\,\%$ bei $h=21$). Damit werden die Risikobänder für mittelfristige Zeiträume zu eng und unterschätzen die Unsicherheit.

Für die Praxis bedeutet dies, dass die Intervalle bei längeren Fristen nicht isoliert interpretiert werden dürfen. Die sinkende Coverage kann dem Einkauf jedoch als Warnsignal dienen, verstärkt auf alternative Absicherungsstrategien wie Lagerhaltung, gestaffelte Beschaffung oder Festpreisvereinbarungen zu setzen.

\subsection{Grenzen der Datenbasis}
Die Beschränkung auf historische Preis- und Volumenzeitreihen sowie gesamt-wirtschaftliche Indizes lässt weitere Faktoren außer Acht.
Variablen wie physische Lagerbestände, Energiepreise als Primärkostenfaktor der Produktion sowie politische Entscheidungen fließen nur indirekt über die Preisbildung in das Modell ein.

Eine Erweiterung der Datenbasis um solche Zeitreihen oder unstrukturierte Daten, etwa durch Text-Mining von Marktberichten, könnte die Signalqualität verbessern, würde jedoch die grundlegenden Grenzen der Vorhersagbarkeit nicht vollständig aufheben.

In der aktuellen Form kann das System daher keine belastbare Kaufempfehlung für einen einzelnen optimalen Zeitpunkt liefern.
Für die betriebliche Praxis sind stattdessen mittel- bis langfristige Beschaffungsstrategien, ein strukturiertes Bestandsmanagement und flexible Liefervereinbarungen die sinnvolleren Instrumente, um Preisrisiken zu begrenzen.
